\documentclass[11pt]{res}

\usepackage[colorlinks,urlcolor={blue}]{hyperref}
\usepackage{charter}
\usepackage{etaremune}

\newenvironment{cvitemize}{\begin{itemize}\itemsep -2pt}{\end{itemize}\vspace{-5pt}}
\newenvironment{cvenumerate}{\begin{etaremune}\itemsep -2pt}{\end{etaremune}\vspace{-5pt}}

\begin{document}
\begin{resume}

\moveleft.5\hoffset
\leftline{\hspace{-18pt}\large\bf Shukai Du\hfill Curriculum Vitae\hspace{-18pt}}
\moveleft\hoffset
\vbox{\hrule width\resumewidth height 1.5pt}
\begin{minipage}[t]{0.5\textwidth}
Assistant Professor\\
Department of Mathematics\\
Syracuse University
\end{minipage}
\begin{minipage}[t]{0.05\textwidth}
\end{minipage}
\begin{minipage}[t]{0.4\textwidth}
Office: Carnegie 206D\\
Email: \href{mailto:sdu113@syr.edu}{sdu113@syr.edu}\\
Website: \href{https://shukaidu.github.io}{https://shukaidu.github.io}
\end{minipage}

\section{RESEARCH INTERESTS}
\begin{cvitemize}
\item
Finite element and discontinuous Galerkin methods
\item
Scientific machine learning and data-driven approaches
\item
Symplectic and structure-preserving methods
\item
Adaptive mesh solver for radiative transfer
\item
Electromagnetic and elastic/viscoelastic waves
\item
Computational inverse and ill-posed problems
\end{cvitemize}

\section{PROFESSIONAL EXPERIENCE}
Syracuse University
\begin{cvitemize}
\item Assistant Professor
\hfill Oct 2024 -- now
\end{cvitemize}

University of Wisconsin-Madison
\begin{cvitemize}
\item Visiting Assistant Professor
\hfill Sep 2020 -- Sep 2024
\end{cvitemize}

University of Minnesota-Twin Cities
\begin{cvitemize}
\item Visiting Doctoral Student 
\hfill Sep 2019 -- June 2020
\end{cvitemize}

\section{EDUCATION}
{\bf University of Delaware}
\begin{cvitemize}
\item {\bf Ph.D} in {\bf Applied Mathematics}  \hfill 2020\\
Advisor: Dr. Francisco-Javier Sayas\\
Thesis: Generalized projection-based error analysis of hybridizable discontinuous Galerkin methods
\end{cvitemize}

{\bf Wuhan University}
\begin{cvitemize}
\item  M.S. in Computational Mathematics\hfill 2015
\item  B.S. in Pure Mathematics\hfill 2012
\end{cvitemize}



\section{PUBLICATIONS}
\begin{cvenumerate}

\item[]
\hspace{-1cm}{\bf Submitted}

\item {B.~Cockburn, {\bf S.~Du}, M. ~A.~Sánchez}. A spectral analysis of hp-hybridized mixed and HDG methods for parametric second-order elliptic problems.

\item {\bf S.~Du}, T.~Le, S.~N.~Stechmann. Element learning with hp-adaptivity: Machine learning-based acceleration of hp-adaptive spectral element methods, and application to atmospheric radiative transfer

\item[]
\item[]

\hspace{-1cm}{\bf Peer-reviewed}
\item
J.~L.~Torchinsky, 
{\bf S. Du}, and S.~N.~Stechmann.
Angular-spatial hp-adaptivity for radiative transfer with discontinuous Galerkin spectral element methods.
{\sl J. Quant. Spectrosc. Radiat. Transf. 348 (2026)}.\\
\href{https://doi.org/10.1016/j.jqsrt.2025.109687}{DOI: 10.1016/j.jqsrt.2025.109687}

\item 
D.~H.~Marsico, {\bf S.~Du}, S.~N.~Stechmann.
Can second-order numerical accuracy be achieved for Moist Atmospheric Dynamics with non-smoothness at cloud edge?
{\sl J. Adv. Model. Earth Syst. 17 (2025).}\\
\href{https://doi.org/10.1029/2025MS005293}{DOI: 10.1029/2025MS005293}

\item
{\bf S. Du}, and S. N. Stechmann.
Element learning: a systematic approach of accelerating finite
element-type methods via machine learning, with applications to radiative transfer. {\sl J. Comp. Math. (2024)}.\\
\href{https://doi.org/10.4208/jcm.2407-m2024-0047}{DOI: 10.4208/jcm.2407-m2024-0047}

\item
{\bf S. Du}, and S. N. Stechmann. Inverse radiative transfer with goal-oriented hp-adaptive mesh refinement: adaptive-mesh inversion. {\sl Inverse Probl. 39 (2023), no. 11}.\\
\href{https://doi.org/10.1088/1361-6420/acf785}{DOI: 10.1088/1361-6420/acf785}

\item
B. Cockburn, {\bf S. Du}, M. A. Sánchez. A priori error analysis of new semidiscrete, Hamiltonian HDG methods for the time-dependent Maxwell's equations.
{\sl ESAIM: M2AN 57 (2023), no.4, 2097-2129}.\\
\href{https://doi.org/10.1051/m2an/2023048}{DOI: 10.1051/m2an/2023048}

\item 
{\bf S. Du}, and S. N. Stechmann. Fast, low-memory numerical methods for radiative transfer via hp-adaptive mesh refinement. {\sl J. Comput. Phys. 480 (2023)}.\\
\href{https://doi.org/10.1016/j.jcp.2023.112021}{DOI: 10.1016/j.jcp.2023.112021}

\item
{\bf S. Du}, and S. N. Stechmann. A universal predictor-corrector approach for minimizing artifacts due to mesh refinement. {\sl J. Adv. Model. Earth Syst. 15 (2023)}.\\
\href{https://doi.org/10.1029/2023MS003688}{DOI: 10.1029/2023MS003688}

\item
B. Cockburn, {\bf S. Du}, M. A. Sánchez. Combining finite element space-discretization with symplectic time-marching schemes for linear hamiltonian systems.
{\sl Front. Appl. Math. Stat. 9 (2023)}.\\
\href{https://doi.org/10.3389/fams.2023.1165371}{DOI: 10.3389/fams.2023.1165371}


\item M. A. S\'{a}nchez, {\bf S. Du}, B. Cockburn, N.-C. Nguyen, J. Peraire. Symplectic Hamiltonian finite element methods for electromagnetics. {\sl Comput. Methods Appl. Mech. Engrg. 396 (2022)}.\\
\href{https://doi.org/10.1016/j.cma.2022.114969}{DOI: 10.1016/j.cma.2022.114969}

\item B. Cockburn, M. A. S\'{a}nchez, {\bf S. Du}. Discontinuous Galerkin methods with time-operators in their numerical traces for time-dependent electromagnetics. {\sl Comput. Meth. Appl. Math. (2022)}.\\
\href{https://doi.org/10.1515/cmam-2021-0215}{DOI: 10.1515/cmam-2021-0215}


\item
{\bf S. Du}, and F.-J. Sayas. A note on devising HDG+ projections on polyhedral elements. {\sl Math. Comp. 90 (2021), 65-79}.\\
\href{https://doi.org/10.1090/mcom/3573}{DOI: 10.1090/mcom/3573}


\item 
{\bf S. Du}. HDG methods for Stokes equation based on strong symmetric stress formulations. {\sl J. Sci. Comput. 85, 8 (2020)}.\\
\href{https://doi.org/10.1007/s10915-020-01309-7}{DOI: 10.1007/s10915-020-01309-7}

\item
{\bf S. Du}, and F.-J. Sayas. A unified error analysis of hybridizable discontinuous Galerkin methods for the static Maxwell equations. {\sl SIAM J. Numer. Anal. 58 (2020), no. 2, 1367–1391}.\\
\href{https://doi.org/10.1137/19M1290966}{DOI: 10.1137/19M1290966}

\item {\bf S. Du}, and F.-J. Sayas. New analytical tools for HDG in elasticity, with applications to elastodynamics. {\sl Math. Comp. 89 (2020), 1745-1782}.\\
\href{https://doi.org/10.1090/mcom/3499}{DOI: 10.1090/mcom/3499}

\item
{\bf S. Du}, and N. Du. A factorization of least-squares projection schemes for ill-posed problems. {\sl Comput. Meth. Appl. Math. 20 (2020), no. 4, 783-798}.\\
\href{https://doi.org/10.1515/cmam-2019-0173}{DOI: 10.1515/cmam-2019-0173}

\item T.S. Brown, {\bf S. Du}, H. Eruslu, and F.-J. Sayas. Analysis of models for viscoelastic wave propagation. {\sl Appl. Math. Nonlin. Sci. 3 (2018), no. 1, 55-96}.\\
\href{https://doi.org/10.21042/AMNS.2018.1.00006}{DOI: 10.21042/AMNS.2018.1.00006}


\item[]
\item[]
\hspace{-1cm}{\bf Books}

\item
{\bf S. Du}, and F.-J. Sayas. An invitation to the theory of the Hybridizable Discontinuous Galerkin Method. {\sl SpringerBriefs in Mathematics (2019)}.\\
\href{https://doi.org/10.1007/978-3-030-27230-2}{DOI: 10.1007/978-3-030-27230-2}
\end{cvenumerate}

%\noindent{\bf Submitted articles}
%\begin{cvenumerate}
%\item
%\end{cvenumerate}

%\noindent{\bf In preparation}
%\begin{cvenumerate}
%\item 
%\end{cvenumerate}

\section{GRANTS}
\begin{cvitemize}
\item NSF (DMS–2324368): Breaking the 1D Barrier in Radiative Transfer: Fast, Low-Memory Numerical Methods for Enabling Inverse Problems and Machine Learning Emulators. Senior personnel. \$498,832 total, \$350,000 at UW (2023–2026).
\item
NSF (AGS–2326631): Convective Processes in the Tropics Across Scales. Senior personnel. \$768,471 total, \$471,155 at UW (2024–2026).
\end{cvitemize}


\section{PRESENTATIONS}
\begin{cvenumerate}
\item[]
\hspace{-1cm}{\bf Invited talks}

\item 
Inverse radiative transfer via goal-oriented adaptive
mesh refinement\\
{
\sl SIAM New York--New Jersey--Pennsylvania Section
\hfill Nov 2025
}

\item
Accelerating finite element-type methods with machine
learning on substructures\\
{\sl Third HKSIAM Biennial Conference, Chinese University of Hong Kong}
\hfill July 2025

\item
Element learning: accelerating finite element-type
methods via machine learning, with applications to
radiative transfer\\
{\sl SIAM sectional meeting, Fort Worth, Texas}
\hfill March 2025

\item
Forward and inverse computation for radiative transfer
via hp-adaptive mesh refinements\\
14th AIMS Conference, Abu Dhabi
\hfill Dec 2024

\item
Inverse radiative transfer via goal-oriented hp-adaptive mesh refinement\\
14th AIMS Conference, Abu Dhabi
\hfill Dec 2024


\item Element learning: accelerating finite element-type methods via operator learning with cost-effective training\\
{\sl SIAM sectional meeting, Rochester Institute of Technology}
\hfill Nov 2024

\item Forward and inverse computation for radiative transfer via hp-adaptive mesh refinement and machine learning acceleration\\
{\sl Analysis seminar, Binghamton University}\hfill Oct 2024

\item Element learning: a systematic approach of accelerating finite element-type methods via machine learning\\
{\sl CCAM seminar, Purdue University}
\hfill Mar 2024

\item
Element learning: a systematic approach of accelerating
finite element-type methods via machine learning\\
{\sl Analysis and Data Science Seminar, SUNY at Albany}
\hfill Feb 2024

\item
Element learning: a systematic approach of accelerating
finite element-type methods via machine learning\\
{\sl Math Department Colloquium, Syracuse University}
\hfill Jan 2024

\item
Element learning: a systematic approach of accelerating
finite element-type methods via machine learning\\
{\sl Math Department Colloquium, Chinese University of Hong Kong}
\hfill Dec 2023

\item 
Element learning: a systematic approach of accelerating finite element-type methods, with applications to radiative transfer\\
{\sl University of Electronic Science and Technology of China}
\hfill Nov 2023

\item 
Element learning: a systematic approach of accelerating finite element-type methods via machine learning, with applications to radiative transfer\\
{\sl Scientific Computing Seminars, University of Houston}
\hfill Nov 2023

\item 
Element learning: a systematic approach of accelerating finite element-type methods via machine learning, with applications to radiative transfer\\
{\sl Applied Math seminar, University of Louisiana at Lafayette}
\hfill Oct 2023

\item 
Element learning: a systematic approach of accelerating finite element-type methods, with applications to radiative transfer\\
{\sl Numerical analysis and PDE seminar, University of Delaware}
\hfill Sep 2023

\item
Energy-conserving discontinuous Galerkin methods
with time-operators in their traces for time-dependent
electromagnetics\\
{\sl 17th UCNCCM, Albuquerque, NM}
\hfill July 2023

\item
Fast, low-memory methods for radiative transfer
through hp-adaptive mesh refinement\\
{\sl 13th AIMS Conference, Wilmington, NC}
\hfill June 2023

\item
Unified analysis of HDG methods for the static Maxwell equations\\
{\sl CILAMCE-PANACM 2021, Brazil} \hfill Nov 2021

\item Generalized projection-based error analysis of
hybridizable discontinuous Galerkin (HDG) methods\\
{\sl CEDYA2021, Spain} \hfill June 2021

\item Unified analysis of HDG methods for the static Maxwell equations\\
{\sl SIAM CSE2021, Virtual Meeting}\hfill Mar 2021

\item
Projection-based analysis of hybridizable discontinuous
Galerkin (HDG) methods\\
{\sl Wenbo Li Prize Talk, U of Delaware}\hfill Feb 2020

\item New analysis techniques of HDG+ method\\
{\sl SIAM Sectional Meeting, Iowa State U}\hfill Oct 2019

\item
Uniform-in-time optimal convergent HDG method for\\
transient elastic waves with strong symmetric stress formulation\\
{\sl WAVES2019, TU Wien, Vienna}\hfill Aug 2019

\item
Hybridizable Discontinuous Galerkin schemes for elastic waves\\
{\sl ICIAM2019, Valencia}\hfill July 2019

\item
HDG for transient elastic waves\\
{\sl WONAPDE2019, U of Concepcion}\hfill Jan 2019

\item[]\item[]
\hspace{-1cm}{\bf Contributed talks}
\item Element learning: accelerating finite element methods via operator learning\\
FEM Circus, U of Notre Dame\hfill Oct 2023

\item 
Three-dimensional radiative transfer: fast, low-memory
numerical methods\\
{\sl Collective Madison Meeting, Madison, WI}
\hfill Aug 2022

\item
Projection-based analysis of HDG methods with reduced stabilization\\
{\sl DelMar Num Day 2019, U of Maryland}\hfill May 2019
\item
Projection-based error analysis of HDG methods for transient elastic waves\\
{\sl FEM Circus, U of Delaware} \hfill Nov 2018
\item
Devising a tailored projection for a new HDG method in linear elasticity\\
{\sl FEM Circus, U of Tennessee} \hfill Mar 2018
\item
A new HDG projection and its applications\\
{\sl Mid-Atlantic Numerical Analysis Day, Temple U} \hfill Nov 2017

\item[]\item[]
\hspace{-1cm}{\bf Poster presentation}

\item Fast, low-memory numerical methods for radiative transfer: forward and inverse problems\\
{\sl New Trends in Computational and Data Sciences, Caltech}\hfill Dec 2022
\item
Hybridizable Discontinuous Galerkin methods
in transient elastodynamics\\
{\sl FACM2018, New Jersey Institute of Technology}\hfill Aug 2018
\item
Building a computational code for 3D viscoelastic wave simulation\\
{\sl Mid-Atlantic Numerical Analysis Day, Temple U} \hfill Nov 2016
\end{cvenumerate}

\section{TEACHING}
{\bf Instructor}
\begin{cvitemize}
\item Methods of Numerical Analysis II (MAT684) at Syracuse University
\hfill Spring 2026
\item Applied Linear Algebra (MAT532) at Syracuse University
\hfill Fall 2025
\item Multivariable calculus (MAT397) at Syracuse University \hfill Spring 2025
\item Linear Algebra and Differential Equations (Math320) at UW-Madison\hfill Spring 2023
\end{cvitemize}

{\bf Teaching Assistant}
\begin{cvitemize}
\item Analytic Geometry and Calculus C (Math243)\hfill Fall 2016 \& Fall 2017
\item Analytic Geometry and Calculus B (Math242)\hfill Spring 2017
\item Calculus I (Math221)\hfill Spring 2018
\item Review of Advanced Mathematical Problems\\
(summer courses offered to incoming graduate students)\hfill Fall 2018
\end{cvitemize}

% {\bf International Teaching Assistant (ITA) training program}
% \begin{cvitemize}
% \item
% Graduated with the highest category of scores (category I) \hfill Summer 2015
% \end{cvitemize}



\section{MENTORING ACTIVITIES}
{\bf Graduate mentorship}
\begin{cvitemize}
\item
Jason Torchinsky (co-mentored with S.~N.~Stechmann)
\hfill 2022 -- 2023
\end{cvitemize}

{\bf Undergraduate mentorship}
\begin{cvitemize}
\item WISCERS project at the University of Wisconsin-Madison
\hfill 2023
% -- a research-focused mentorship program for undergraduate students

\item 
GEMS summer research project at the University of Delaware
\hfill Fall 2016
\end{cvitemize}

\section{JOURNAL REFEREE}
Journal of Scientific Computing \textbullet\
SIAM Journal on Scientific Computing \textbullet\
SIAM Multiscale Modelling and Simulation \textbullet\
ESAIM: Mathematical Modelling and Numerical Analysis \textbullet\
Inverse Problems and Imaging \textbullet\
Computers and Mathematics with Applications \textbullet\
Frontiers in Applied Mathematics and Statistics

\section{AWARDS AND HONORS}
Wenbo Li Prize \hfill 2020\\
University Doctoral Fellowship Award at the University of Delaware \hfill 2019\\
ICIAM2019 travel grant\hfill 2019 \\
Graduate Enrichment Fellowship at the University of Delaware\hfill 2018\\
GEMS project fund at the University of Delaware\hfill Summer 2016\\
National Scholarship for Graduate Students of China\hfill 2013\\
People's Scholarship of Wuhan University\hfill 2011\\
Outstanding Student of Wuhan University\hfill 2009 -- 2011

\section{SERVICES}
\begin{itemize}
\item Organizing minisymposium ``Recent Developments on Hybrid Methods Combining Neural Networks with Classical Numerical Methods"\\
SIAM Pacific Northwest Section \hfill Oct 2025
\item Mentor for SIAM Virtual Resume Building Workshop\hfill Oct 2024
\end{itemize}



\section{CODING PROJECTS}
Fast, low-memory methods for radiative transfer
\hfill 2020 -- 2022
\begin{itemize}
\item Built a cell-based structured adaptive mesh refinement (AMR) data structure
\item Implemented discontinuous Galerkin (DG) methods with $hp$-adaptivity for the full radiative transfer equation
\end{itemize}

Hybridizable Discontinuous Galerkin (HDG) methods \hfill 2016 -- 2020\\
(based on 
\href{https://github.com/team-pancho/HDG3D}{HDG3D library}:
github.com/team-pancho/HDG3D) 
\begin{cvitemize}
\item Built Matlab codes of high order HDG methods on computing cluster for transient elastic/viscoelastic waves and Maxwell equations
\item Wrote documentation with detailed implementation procedures for HDG methods for Maxwell equations
\end{cvitemize}

Finite Element Method (FEM) \hfill 2016\\
(based on 
\href{https://team-pancho.github.io}{Team Pancho}
FEM library:
team-pancho.github.io) 
\begin{cvitemize}
\item Built Matlab codes of high order FEM methods on computing cluster for simulation of viscoelastic waves.
\end{cvitemize}

Multiscale modeling \hfill 2013 -- 2015
\begin{cvitemize}
\item Implemented algorithms to calculate Cauchy stress tensor based on micro-scale molecular dynamics information

\end{cvitemize}

\section{COMPUTER SKILLS}
{\sl Theory}\\
Data Structures $\bullet$ Algorithm $\bullet$
Object Oriented Programming\\
{\sl Languages \& Software}\\
Matlab $\bullet$ Python $\bullet$ C $\bullet$ C++ $\bullet$ Fortran $\bullet$ 
openMPI $\bullet$ LISP $\bullet$ Linux Shell


%\section{ACTIVITIES}
%MSRI Summer School on Harmonic Analysis, Park City
%\hfill Jul 2018\\
%Nonlocal School on Fractional Equations, Iowa State U
%\hfill Aug 2017\\
%Finite Element Circus, Rutgers U \hfill Apr 2017\\
%Summer School on Applied Mathematics in Beijing University \hfill Jul 2014\\
%Second Pacific Rim Mathematical Association Congress \hfill Jun 2013\\
%International Conference on Mathematical Modeling \& Computation \hfill May 2013\\
%Summer School on Statistical Learning and Inference for \\ Massive Data in Fudan University \hfill Jul 2012

\let\thefootnote\relax\footnotetext{Last update: \today}


\end{resume}
\end{document}